\section{Introduction}

Electroencephalography (EEG) is cheap, portable and widely available, which makes it an attractive candidate for
objective markers of major depressive disorder (MDD). Many machine-learning studies report high accuracies for
separating patients with MDD from healthy controls (HC) using EEG. However, a large share of this literature evaluates
models by splitting \emph{trials} or \emph{segments}, rather than \emph{participants}, between training and test sets.

When the same person contributes data to both sets, a classifier can succeed by recognising the individual (their
head geometry, impedances, idiosyncratic rhythms or recording session) instead of the disorder. The resulting estimates
do not approximate the intended use-case, which is diagnosing a new person~\cite{saeb2017need}. This form of leakage
is a documented driver of the reproducibility crisis in ML-based science~\cite{kapoor2023leakage}. It has recently been
shown to inflate performance substantially in translational EEG studies~\cite{brookshire2024leakage}.

Small samples aggravate the problem. With a few dozen participants, cross-validated accuracy has error bars of
$\pm10$ percentage points or more~\cite{varoquaux2018cross}, and theoretical chance levels are not attained in
practice~\cite{combrisson2015exceeding}. Finally, demographic or acquisition variables that differ between groups
(sex, education, recording period) can masquerade as disorder-related signal.

The Multi-modal Open Dataset for Mental-disorder Analysis (MODMA)~\cite{cai2022modma} provides 128-channel EEG
recorded during an emotional-face dot-probe task~\cite{macleod1986attentional} from patients with MDD and matched
controls. It is a natural testbed for these questions. In this paper we ask a single question:
\emph{how well can task EEG classify MDD versus HC in participants that the model has never seen?}

\paragraph{Contributions.}
\begin{enumerate}
  \item A complete pipeline from the raw EGI files, with exact ID matching against the clinical spreadsheet,
        event decoding by channel name, per-participant preprocessing (PREP bad-channel detection, ICA with ICLabel)
        and pre-specified participant exclusion criteria.
  \item An analysis plan committed to version control before any result existed. It fixes the primary model, the
        primary outcome and all thresholds, and every deviation is reported.
  \item A strictly subject-wise evaluation of 18 model configurations across five families: a demographic baseline, handcrafted
        biomarkers (spectral, frontal alpha asymmetry, ERP, connectivity, complexity), Riemannian geometry and three
        deep networks.
        \begin{itemize}
          \item Every model uses nested cross-validation and identical subject partitions.
          \item Every result carries a participant-level bootstrap confidence interval; all models except the
                deep networks also carry a label-permutation test of the full nested pipeline.
        \end{itemize}
  \item Confound analyses: can EEG predict sex, age or recording period, and does performance survive restriction
        to a period in which both groups were recorded?
  \item A direct measurement of how much trial-wise splitting inflates performance on the same data.
\end{enumerate}
